\documentclass[aspectratio=169, 12pt]{beamer}

% --- Language and Fonts ---
\usepackage{fontspec}
\usepackage[english]{babel}
\usepackage{amsmath, amssymb, bm}
% tikz
\usepackage{tikz}
\usetikzlibrary{shapes, arrows.meta, positioning, shadows, calc}

% --- Bibliography (BibTeX + natbib) ---
\usepackage{natbib}
\bibliographystyle{chicago} 

% =========================================
% ===== keybox ====================
% =========================================
\usepackage[most]{tcolorbox}
\newtcolorbox{keybox}[1]{
  enhanced,
  colback=headbg!4,
  colframe=headbg,
  coltitle=headfg,
  colbacktitle=headbg,
  title=#1,
  fonttitle=\bfseries,
  boxrule=0.8pt,
  arc=1.5mm,
  left=2mm,
  right=2mm,
  top=1mm,
  bottom=1mm,
  before skip=0.2cm,
  after skip=0.2cm
}

\usepackage{booktabs}

% =========================================
% ===== Custom Footline (Page Number) =====
% =========================================


\setbeamercolor{footline}{fg=headbg}
\setbeamerfont{footline}{size=\scriptsize}

\setbeamertemplate{footline}{%
  \hfill%
  \usebeamercolor[fg]{footline}%
  \usebeamerfont{footline}%
  \insertframenumber\,/\,\inserttotalframenumber\kern1em\vskip3pt%
}

\setbeamertemplate{navigation symbols}{}

\usetheme{default}
\usepackage{xcolor}

\definecolor{headbg}{RGB}{30, 60, 120}
\definecolor{headfg}{RGB}{255, 255, 255}
\definecolor{accent}{RGB}{200, 50, 50}

\setbeamercolor{headline}{bg=headbg, fg=headfg}
\setbeamercolor{section in head/foot}{bg=headbg, fg=headfg}
\setbeamercolor{subsection in head/foot}{bg=headbg!80, fg=headfg}
\setbeamercolor{frametitle}{bg=headbg, fg=headfg}

% Step indicator command
\newcommand{\stepindicator}[1]{%
  \hfill {\scriptsize \textit{\textcolor{headbg}{Step #1 / 6}}}%
}
\DeclareMathOperator*{\argmin}{arg\,min}

\usepackage{hyperref}


% --- Title ---
\title[]{A First Course in Causal Inference Ch.~10}
\author[Kodai Minesaki]{Kodai Minesaki}
\institute[Faculty of Engineering, The University of Tokyo]{Faculty of Engineering, The University of Tokyo}
\date{\today}

\begin{document}

% --- Title Slide ---
\begin{frame}
\titlepage
\end{frame}


\begin{frame}

  {
    \large
    Ch.~10 Observational Studies, Selection Bias, and Nonparametric Identification of Causal Effects
  }
  
\end{frame}

\begin{frame}{Ch.~10 Observational Studies, Selection Bias, and Nonparametric Identification of Causal Effects}
  \citet{cochran1965planning} summarized two commonn characteristics of observational studies:
  \begin{enumerate}
    \item the objective is to elucidate cause-and-effect relationships;
    \item it is not feasible to use controlled experimentation.
  \end{enumerate}
  \begin{itemize}
    \item This part will discuss causal inference with observational studies.
    \item Many ideas of causal inference with observational studies are deeply connected to those with randomized experiments.
  \end{itemize}
\end{frame}

\begin{frame}{10.2 Causal effects and selection bias under the potential outcomes framework}
  \begin{itemize}
    \item For unit $i$ $(i=1,\dots,n)$, we have pretreatment covariates $X_i$, a binary treatment indicator $Z_i$, and an observed outcome $Y$ with two potential outcomes $Y_i(1)$ and $Y_i(0)$.
    \item We assume $\{X_i,Z_i,Y_i(1),Y_i(0)\}_{i=1}^n\overset{\text{i.i.d.}}{\sim}\{X,Z,Y(1),Y(0)\}.$
    \item The causal effects of interest are the average causal effect $\tau=\mathbb{E}[Y(1)-Y(0)]$, the average causal effect on the treated units $\tau_T=\mathbb{E}[Y(1)-Y(0)\mid Z=1]$, and the average causal effect on the control units $\tau_C=\mathbb{E}[Y(1)-Y(0)\mid Z=0].$
  \end{itemize}
\end{frame}

\begin{frame}{10.2 Causal effects and selection bias under the potential outcomes framework}
  Rewrite the average causal effects on the treated units and control units:
  \[
  \tau_T=\mathbb{E}[Y(1)\mid Z=1]-\mathbb{E}[Y(0)\mid Z=1]=\underbrace{\mathbb{E}[Y\mid Z=1]}_{\text{observable}}-\underbrace{\mathbb{E}[Y(0)\mid Z=1]}_{\text{unobservable}}.
  \]
  \[
  \tau_C=\mathbb{E}[Y(1)\mid Z=0]-\mathbb{E}[Y(0)\mid Z=0]=\underbrace{\mathbb{E}[Y(1)\mid Z=0]}_{\text{unobservable}}-\underbrace{\mathbb{E}[Y\mid Z=0]}_{\text{observable}}.
  \]
  These unobservable terms are \textit{counterfactuals}.
\end{frame}


\begin{frame}{10.2 Causal effects and selection bias under the potential outcomes framework}
  \begin{itemize}
    \item The simple difference in means, known as the \textit{prima fucie} causal effect
    \[
    \tau_{\text{PF}}=\mathbb{E}[Y\mid Z=1]-\mathbb{E}[Y\mid Z=0]=\mathbb{E}[Y(1)\mid Z=1]-\mathbb{E}[Y(0)\mid Z=0]
    \]
    is generally biased for the causal effects defined above.
    \item For example, $\tau_{\text{PF}}-\tau_T=\mathbb{E}[Y(0)\mid Z=1]-\mathbb{E}[Y(0)\mid Z=0]$ is not zero in general.
    \item The equation above quantifies \textit{selection bias}.
    \item The reason why randomization matters is because the \textit{selection bias} term is zero.
    \item Under the CRE, $\tau_{\text{PF}}=\tau_T=\tau_C=\tau$.
    \item The fundamental benefit of randomization is to balance the distributions of the potential outcomes across the groups.
  \end{itemize}
\end{frame}


\begin{frame}{10.3 Sufficient conditions for nonparametric identification of causal effects}
  \begin{itemize}
    \item Causal inference with observational studies is challenging and it relis on strong assumptions.
    \item A strategy is to use the information from the pretreatment covariates and assume that conditioning on observed covariates $X$, the selection bias terms are zero.
    {\small\[
    \mathbb{E}[Y(0)\mid Z=1,X]=\mathbb{E}[Y(0)\mid Z=0,X],\quad \mathbb{E}[Y(1)\mid Z=1,X]=\mathbb{E}[Y(1)\mid Z=0,X]
    \]}
    \item This assumption states that the differences in the means of the potential outcomes across the treatment and control groups are entirely due to the difference in the observed covariates.
    \item The average causal effect $\tau$ is nonparametrically identifiable under this assumption.
  \end{itemize}
\end{frame}


\begin{frame}{10.3 Sufficient conditions for nonparametric identification of causal effects}
  \begin{keybox}{Definition~10.1 (identification)}
    A parameter $\theta$ is identifiable if it can be writte as a function of the distribution of the observed data under certain model assumptions. A parameter $\theta$ is nonparametrically identifiable if it can be written as a function of the distribution of the observed data without any parametric model assumptions.
  \end{keybox}
  \begin{itemize}
    \item This assumption is fundamental in causal inference with observational studies.
    \item The parameter of interest $\tau=\mathbb{E}[Y(1)-Y(0)]$ depends on some unobserved random variables, so it is unclear whether it is nonparametrically identifiable based on observed data.
    \item Because $\tau_{\text{PF}}$ depends only on the observables, it is nonparametrically identified by definition.
  \end{itemize}
\end{frame}

\begin{frame}{10.3 Sufficient conditions for nonparametric identification of causal effects}
  \begin{keybox}{Theorem~10.1}
    Under the assumption, the average causal effect $\tau$ is identified by
    \begin{align*}
      \tau &=\mathbb{E}[\tau(X)]\\
      &= \mathbb{E}[\mathbb{E}[Y\mid Z=1,X]]-\mathbb{E}[Y\mid Z=0,X]\\
      &= \int \{\mathbb{E}[Y\mid Z=1,X=x]-\mathbb{E}[Y\mid Z=0,X]\}f(x)dx
    \end{align*}
  \end{keybox}
\end{frame}

\begin{frame}{10.3 Sufficient conditions for nonparametric identification of causal effects}
  Usually we impose the strong assumption:
  \begin{keybox}{Assumption~10.1 (ignorability)}
    We have $Y(z)\perp\!\!\!\perp Z\mid X\quad (z=0,1)$.
  \end{keybox}
  \begin{keybox}{Assumption~10.2 (strong ignorability)}
    We have $\{Y(1),Y(0)\}\perp\!\!\!\perp  Z \mid X$.
  \end{keybox}
  These assumptions are not necesarry for identifying $\tau$ but they cannot be relaced if the parameter of interest is the causal effects on other scales.
\end{frame}

\begin{frame}{10.3 Sufficient conditions for nonparametric identification of causal effects}
  \begin{itemize}
    \item Intuitively, the assumption rules out all unmeasured covariates that affect the treatment and outcome simultaneously.
    \item The ignorability assumption can be reasonable if we observe a rich set of covariates $X$ that affect the treatment and the outcome simultaneously.
    \item It is fundamentally untestable.
    \item We may justify it based on the scientific background knowledge but we are often not sure whether it holds or not.
  \end{itemize}
\end{frame}

\begin{frame}{10.4 Two simple estimation strategies and their limitations}
  \begin{itemize}
    \item If the covariate $X_i\in\{1,\dots,K\}$ is discrete then the ignorability assumption reads as $Y(z)\perp\!\!\!\perp Z\mid X=k$. 
    \item Therefore, we can use the estimator $\hat{\tau}=\sum_{k=1}^{K}\pi_{[k]}\{\hat{\bar{Y}}(1)-\hat{\bar{Y}}(0)\},$ where $\pi_{[k]}$ denotes a propensity score.
    \item This method is widely used but there are several problems: 
    \item First, it works well for small $K$, but it leads to the poorly defines $\hat{\tau}_{[k]}$ for large $K$ because some strata are very likely to have $n{[k]1}=0$ or $n_{[k]1}=0$.
    \item Second, it is not obvious how to apply this stratification method to multidimensional continuous or mixed covariates $X$.
  \end{itemize}
\end{frame}

\begin{frame}{10.4 Two simple estimation strategies and their limitations}
  \begin{itemize}
    \item The most commonly used method is to run the OLS with an additive model of the observed outcome on the treatment indicator and covariates: $\mathbb{E}[Y\mid Z,X]=\beta_0+\beta_zZ+\beta_x^{\top}X$.
    \item If the linear model is correct, then we have \[\tau(X)=\mathbb{E}[Y\mid Z=1,X]-\mathbb{E}[Y\mid Z=0,X]=\beta_z.\]
    \item This estimator implies that the causal effect is homogeneous with respect to the covariates.
    \item Thie, coupled with ignorability, implies that $\tau=\mathbb{E}[\tau(X)]=\beta_z$.
    \item If ignorability holdes and the outcome model is linear, then the average causal effect equals the coeff. of $Z$.
    \item At the same time, this is valid only under two strong assumptions: \textit{ignorability} and \textit{the linear model}.
  \end{itemize}
\end{frame}

\begin{frame}{10.4 Two simple estimation strategies and their limitations}
  \begin{itemize}
    \item If we incorporate heterogenous causal effects, we include theinteraction term:
    \[
    \mathbb{E}[Y\mid Z,X] = \beta_0+\beta_zZ+\beta_x^{\top}X+\beta_{zx}^{\top}XZ.
    \]
    \item We then have
    \[
    \tau(X) = \mathbb{E}[Y\mid Z=1,X] - \mathbb{E}[Y\mid Z=0,X] = \beta_z+\beta_{zx}^{\top}X.
    \]
    \item With ignorability, it implies that 
    \[
    \tau = \mathbb{E}[\tau(X)] = \mathbb{E}[\beta_z+\beta_{zx}^{\top}X] = \beta_z + \beta_{zx}^{\top}\mathbb{E}[X].
    \]
    \item The estimator for $\tau$ is then $\hat{\beta}_z + \hat{\beta}_{zx}^{\top}\bar{X}.$
  \end{itemize}
\end{frame}

\begin{frame}{10.4 Two simple estimation strategies and their limitations}
  \begin{itemize}
    \item If we center the covariates to ensure $\bar{X}=0$, then the estimator is simpli the regression coeff. of $Z$.
    \item In general, we can use other more complex models to estimate the causal effects.
    \item If we build two predictors $\hat{\mu}_1(X)$ and $\hat{\mu}_0(X)$ based on the treated and control data, respectively. 
    \item Then we have an estimator for CATE: $\hat{\tau}(X)=\hat{\mu}_1(X)-\hat{\mu}_0(X)$, and an estimator for ATE: $\hat{\tau}^{\text{reg}}=n^{-1}\sum_{i=1}^{n}\{\hat{\mu}_1(X)-\hat{\mu}_0(X)\}.$
    \item It is sometimes called the \textit{outcome regression} estimator.
  \end{itemize}
\end{frame}

\begin{frame}{10.4 Two simple estimation strategies and their limitations}
  \begin{itemize}
    \item The biggest problem with the above approach based on outcome regressions is its sensitivity to the specification of the outcome model.
    \item Researchers migh report their favorable causal effects estimates after searching over a wide set of candidate models, without confessing the model searching process.
    \item This is a major source of $p$-hacking in causal inference.
  \end{itemize}
\end{frame}


% =========================

% =========================
\section{}
\begin{frame}[allowframebreaks, noframenumbering]{References}
  \beamerdefaultoverlayspecification{}
  \small
  \bibliography{../references}
\end{frame}

\end{document}